% ============================================================
% Section 1: Introduction
% ============================================================

\section{Introduction}
\label{sec:intro}

Real-world data is rarely accessible in safety-critical industries.
A single P\&ID digitization project at an industrial facility can require
months of expert annotation and significant cost — yet the resulting data
cannot be shared externally due to intellectual property constraints.
Piping and Instrumentation Diagrams (P\&IDs) are the primary engineering
document of process plants in the oil, gas, chemical, and power industries,
encoding every instrument, valve, pump, and pipe connection.
Maintaining a digital, queryable process graph from P\&IDs is critical for
safety audits, maintenance scheduling, and digital-twin
workflows~\cite{toghraei2019}.
Despite this industrial urgency, the only publicly available P\&ID dataset
with graph-level ground truth contains just 12 real-world images~\cite{sturmer2025}.

State-of-the-art end-to-end digitization requires tens of annotated real
P\&IDs for training.
Stürmer~\etal~\cite{sturmer2025} train on 60 real P\&IDs combined with
2,000 template-based synthetic images, achieving 75.46\% edge mAP on
OPEN100 --- but explicitly note that ``the Relationformer's performance
suffers significantly when trained on synthetic data only''
(Appendix~A5).
The standard response to data scarcity --- generate synthetic training
images --- does not resolve the problem when generation is template-based:
randomly placing symbols on a canvas yields graphs that do not reflect
industrial process topology, producing only ${\sim}33\%$ edge mAP on real
benchmarks~\cite{sturmer2025}.

We identify the root cause: the domain gap in P\&IDs is
\emph{structural}, not visual.
A model trained on synthetic graphs where components connect in random,
unrealistic topologies learns the wrong priors for real-world layouts.
Starting from just 12 real P\&IDs as seeds, our topology-preserving
generator creates \ours{} --- 665 training images whose graph statistics
mirror those of real industrial diagrams, with no additional annotation.
We pair this dataset with a patch-based inference strategy that processes
high-resolution P\&IDs at appropriate scale, following the approach of
\cite{sturmer2025} and adapting it for our synth-trained pipeline.

Training Relationformer exclusively on \ours{} achieves \textbf{63.8\%
edge mAP} on OPEN100 with \emph{zero real training data} --- 30
percentage points above template-based synthetic, 18 points above the
modular baseline that uses real data, and within 8 points of our
real-data oracle.
Adding even a small amount of real data (E3~Mixed) pushes to 72.6\% edge
mAP, approaching Stürmer~\etal's~\cite{sturmer2025} fully supervised result
of 75.46\% while requiring far fewer real P\&IDs.

This paper makes three contributions:

\begin{itemize}[noitemsep, topsep=3pt]
    \item \textbf{\ours{}}: the first topology-preserving synthetic P\&ID
          dataset with graph-level annotations, released with the
          generation pipeline and trained models upon acceptance.
    \item \textbf{Empirical evidence that generation quality is the
          bottleneck}: consistent with~\cite{sturmer2025} Appendix~A5,
          template-based synth achieves ${\sim}33\%$ edge mAP;
          our topology-preserving approach raises this to 63.8\%
          (+30~pp) using the same model.
    \item \textbf{A rigorous synthetic-to-real benchmark}: 5-fold
          cross-validation $\times$ 3~seeds on OPEN100 provides a
          reproducible evaluation framework for the data-scarce P\&ID
          digitization regime.
\end{itemize}
